\section{The complement and its integral attachments}\label{sec:attachments}

All coordinates in this section refer to the ring $R$ of
Section~\ref{sec:ring}. In particular the ambient multiplication
and the exceptional classes have already been calculated. We use
excision to determine the complement and its maps, with integral
coefficients throughout.

\subsection{The exceptional ideal}
First work in chamber $A$, with $E=P_0\cup P_2$ and $V=M_A\setminus E$.
The planes meet at one point, so a regular neighborhood retracts
to $\PP^2\vee\PP^2$. Integral Poincar\'e--Lefschetz duality on the regular neighborhood,
followed by excision \cite[Theorem~3.43]{hatcher}, gives
\begin{equation}\label{eq:relative}
 H^k(M_A,V)\simeq H_{8-k}(E).
\end{equation}
The only nonzero relative groups are in degrees $4,6,8$, with
ranks $2,2,1$. Their maps to the absolute cohomology are the
plane, line and point inclusions, respectively.

\begin{proposition}\label{prop:complement}
Restriction is a surjective ring homomorphism with kernel
$(p_0,p_2)$. Hence $H^*(V,\Z)=R/(p_0,p_2)$, with free ranks
$1,23,274,21$ in degrees $0,2,4,6$ and no other groups.
\end{proposition}

\begin{proof}
The witnesses \eqref{eq:witness-restrictions} show that the
relative plane and line maps in \eqref{eq:relative} are
primitive injections. The point map is an isomorphism. The
long exact sequence of the pair therefore gives surjectivity
of restriction, vanishing of odd groups, and the stated free
quotients. It remains to identify its kernel as an ideal,
rather than merely as three additive subgroups.

The projection formula on the two planes gives
\begin{equation}\label{eq:ideal-products}
 p_i a=d_i(a)\ell_i\quad(a\in R^2),\qquad
 p_i z=\left(\int_Yp_i z\right)[\pt]\quad(z\in R^4).
\end{equation}
Each $d_i$ takes the value one, and each plane evaluation
takes the value one, by the same unimodular witnesses. Thus
the ideal contains $\ell_0,\ell_2$ and $[\pt]$. Its degree-four
part is exactly the span of $p_0,p_2$ and it has no further
possible degrees. This is the entire restriction kernel
computed from \eqref{eq:relative}, proving the ring statement.
\end{proof}

For every normalized index, restriction to $V$ is consequently
the class $c_k(\widehat J_z)$ of Proposition~\ref{prop:indices}
reduced modulo this explicit ideal. No complement multiplication
is an additional input to that formula.

\subsection{The two-plane neighborhood and its boundary}
Choose a smoothed closed regular neighborhood $W$ of $E$ and
put $L=\partial W$, with the outward boundary orientation from
the complex orientation of $W$. It is the plumbing of the
cotangent normal disk bundles of the two planes. Their normal
Euler numbers are three and their transverse intersection is
positive one. The middle intersection matrix is therefore
\begin{equation}\label{eq:B}
 B=\begin{pmatrix}3&1\\1&3\end{pmatrix}.
\end{equation}
A deleted collar retracts to $L$, so $L$ models the actual
overlap in the cover of $M_A$ by $V$ and an open enlargement
of $W$. It is a boundary inside the resolution. The singular
surface of the target passes through $o$; $L$ is not the link
of an isolated target singularity.

Let $h_i$ be the positive hyperplane classes of the two
components and $t_i=h_i^2$. The neighborhood ring is the
wedge ring
\begin{equation}\label{eq:W-ring}
 h_0h_2=0,\qquad h_i^2=t_i,
\end{equation}
with all products of degree greater than four zero.
Restriction from $R$ is
\begin{equation}\label{eq:W-restrictions}
 a\longmapsto d_0(a)h_0+d_2(a)h_2,\qquad
 z\longmapsto\left(\int_Yzp_0\right)t_0+
                    \left(\int_Yzp_2\right)t_2
\end{equation}
in degrees two and four. It preserves the unit and is zero
in higher positive degrees.

The Poincar\'e--Lefschetz sequence for $(W,L)$ gives
\begin{equation}\label{eq:L-groups}
 H^k(L,\Z)=
 \begin{cases}
 \Z,&k=0,7,\\
 \Z^2,&k=2,5,\\
 \Z/8,&k=4,\\
 0,&\text{otherwise}.
 \end{cases}
\end{equation}
Indeed its only nontrivial even map in the middle is $B$,
whose determinant is eight and whose entries have gcd one.
The quotient map $\Z^2/B\Z^2\to\Z/8$ is
$(n_0,n_2)\mapsto n_0-3n_2$.
Write $x_i$ for the images of $h_i$ and $\tau$ for the
image of $t_0$. If $\omega$ is the oriented degree-seven
generator, choose $\eta_0,\eta_2\in H^5(L,\Z)$ dual to $x_0,x_2$.
The complete nonzero positive products are
\begin{equation}\label{eq:L-products}
 x_0^2=\tau,\quad x_2^2=-3\tau,\quad
 x_i \eta_j=\delta_{ij}\omega,\qquad 8\tau=0,
\end{equation}
and $x_0x_2=0$. Degree and \eqref{eq:L-groups} make all
other positive products zero. The free pairings follow from
oriented Poincar\'e duality; the torsion products follow by
restricting \eqref{eq:W-ring}. These statements describe the
ordinary ring, without asserting an additional torsion linking form.

\begin{proposition}\label{prop:boundary}
With the preceding choices, restriction $\rho:H^*(V)\to H^*(L)$
is
\begin{equation}\label{eq:rho}
 \rho^2(a)=d_0(a)x_0+d_2(a)x_2,\qquad
 \rho^4([z])=\kappa([z])\tau,\quad
 \kappa([z])=\int_Yz(p_0-3p_2)\pmod8.
\end{equation}
It preserves the unit, and its other positive-degree maps vanish.
For the ordered cover $(V,W)$ with alternating map $r_V-r_W$,
the Mayer--Vietoris connecting homomorphisms are
\begin{equation}\label{eq:MV}
 \partial_{\mathrm{MV}}\eta_i=-\ell_i,\qquad
 \partial_{\mathrm{MV}}\omega=-[\pt].
\end{equation}
\end{proposition}

\begin{proof}
Restrictions from $M_A$ to the overlap agree whether taken
through $V$ or $W$. Surjectivity in
Proposition~\ref{prop:complement} and
\eqref{eq:W-restrictions} therefore give \eqref{eq:rho}.
It is well-defined: evaluation of $p_0-3p_2$ against
$p_0,p_2$ gives $0,-8$. Multiplicativity may also be checked
without referring back to the geometric map. For $a,b\in R^2$,
\[
 \int_Yabp_i=d_i(a)d_i(b),
\]
which is exactly the product prescribed by
\eqref{eq:L-products}. All further positive products have
zero target. This checks the map on the complete integral
basis and every ring relation.

For signs, the pair boundary of $(W,L)$ sends the degree-five
classes dual to $h_i$ to the positive relative line classes,
and its oriented top class to the positive relative point
class, by Stokes' convention with outward boundary. Excision
identifies their images in $M_A$ with $\ell_i,[\pt]$.
In the ordered alternating restriction $r_V-r_W$, a
cochain supported on the $W$ side has the negative of that
pair boundary. This yields \eqref{eq:MV}.
\end{proof}

\begin{figure}[htbp]
\centering
\begin{adjustbox}{max width=\textwidth,center}
\begin{tikzpicture}[>=Stealth,font=\small,
 box/.style={draw,rounded corners=2pt,fill=blue!3,align=center,
             minimum height=12mm,minimum width=31mm},
 label/.style={fill=white,inner sep=2pt,font=\footnotesize}]
\node[box] (M) at (0,0) {$R=H^*(M_A,\Z)$};
\node[box] (W) at (6.2,0) {$H^*(W,\Z)$\\two-plane wedge ring};
\node[box] (V) at (0,-2.1) {$H^*(V,\Z)$\\$R/(p_0,p_2)$};
\node[box] (L) at (6.2,-2.1) {$H^*(L,\Z)$\\$H^4(L)=\Z/8$};
\draw[->,thick] (M) -- node[label,above] {restriction} (W);
\draw[->,thick] (M) -- node[label,left] {quotient} (V);
\draw[->,thick] (W) -- node[label,right] {restriction} (L);
\draw[->,thick] (V) -- node[label,above] {$\rho$} (L);
\node at (3.1,-3.0) {$\rho^4([z])=\int z(p_0-3p_2)\bmod8$};
\node[align=center] at (3.1,-3.9)
 {$\partial_{\rm MV}(\eta_i)=-\ell_i,\qquad
   \partial_{\rm MV}(\omega)=-[\mathrm{pt}]$\\[2pt]
  ordered difference $r_V-r_W$, outward boundary orientation};
\end{tikzpicture}

\end{adjustbox}
\caption{The computed restriction square in chamber $A$.
The lower map in degree four is $\kappa$ and the right map
is the cokernel of $B$. The signed connecting maps belong
to the ordered Mayer--Vietoris sequence and are separate
from these ring homomorphisms.}
\label{fig:excision}
\end{figure}

In degree four the square recovers the additive compatibility
description
\begin{equation}\label{eq:fiber-product}
 R^4\simeq
 \left\{([z],n_0,n_2):\ n_0-3n_2\equiv
          \int_Yz(p_0-3p_2)\pmod8\right\},
\end{equation}
where $[z]\in(R/(p_0,p_2))^4$. The ambient ring was already
constructed in Section~\ref{sec:ring};
\eqref{eq:fiber-product} is not a separate derivation of
multiplication from an additive sequence.

For the normalized indices, \eqref{eq:index-plane} gives
\begin{equation}\label{eq:index-kappa}
 \begin{aligned}
 \kappa(c_2(I_z^y))&=\binom{m_0(z)}2-3\binom{m_2(z)}2\pmod8,\\
 \kappa(c_1(I_z^y)c_1(I_{z'}^y))
     &=m_0(z)m_0(z')-3m_2(z)m_2(z')\pmod8.
 \end{aligned}
\end{equation}
These reproduce the earlier binomial attachment of
\cite[\S6.4]{wuebben17}. The present formula also calculates
its full domain: evaluate on $A,T_i,C_{ij},Q_i$ using
\eqref{eq:gram}, with $p_i$ expanded by
\eqref{eq:degree2-products}. This is a specified integral
generating set with all its products, not a test on a
possibly proper sublattice of degree four.

\subsection{A uniform formula for all six chambers}
In chamber $D$ the two planes have the same surface classes
and the negative line classes. Thus use the preceding formulas
with degree-two restriction $a\mapsto\epsilon d_0(a)h_0+
\epsilon d_2(a)h_2$, where $\epsilon=1$ in $A$ and $-1$ in
$D$. The relative maps in degrees $4,6,8$ are
\begin{equation}\label{eq:two-plane-relative}
 (p_0,p_2),\qquad (\epsilon\ell_0,\epsilon\ell_2),\qquad[\pt].
\end{equation}
The boundary ring and degree-four functional are unchanged,
and the connecting maps are the negatives of the displayed
line and point maps.

For the four ruled configurations use
\begin{equation}\label{eq:ruled-template}
 \begin{array}{c|rrrr}
 \sigma&B&C&E&F\\\hline
 (k,j,\epsilon)&(2,0,-1)&(0,2,1)&(2,0,1)&(0,2,-1).
 \end{array}
\end{equation}
The surfaces are $p_k$ and $p_j-p_k$. Write $e_\Sigma$ for
the negative section and $r_\Sigma$ for the ruling in
$\Sigma\simeq\FF_1$. Their global curve classes are
\begin{equation}\label{eq:ruled-curves}
 [e_\Sigma]=\epsilon c_k,\qquad
 [r_\Sigma]=-\epsilon(c_0+c_2),\qquad
 e_\Sigma^2=-1,\quad e_\Sigma r_\Sigma=1,\quad r_\Sigma^2=0.
\end{equation}
Let $\xi,\upsilon\in H^2(W_\sigma)$ take values $(1,0)$
and $(0,1)$ on $(e_\Sigma,r_\Sigma)$. Their restrictions
to the plane are $h,0$ and to $\Sigma$ are
$r_\Sigma,e_\Sigma+r_\Sigma$, respectively. With
$t_P,t_\Sigma$ dual to the two surface components, the ring is
\begin{equation}\label{eq:ruled-W}
 \xi^2=t_P,\qquad \xi\upsilon=\upsilon^2=t_\Sigma,
\end{equation}
with higher positive products zero. This follows from the
Mayer--Vietoris ring of the two components meeting along
their common curve, or directly from the two restrictions.

For $a\in R^2$ set $s=\epsilon d_k(a)$ and
$t=-\epsilon(d_0(a)+d_2(a))$. The complete positive
restrictions to the neighborhood are
\begin{equation}\label{eq:ruled-restriction}
 \begin{aligned}
 a|_{W_\sigma}&=s\xi+t\upsilon,&
 a|_\Sigma&=t e_\Sigma+(s+t)r_\Sigma,\\
 z|_{W_\sigma}&=\left(\int_Yzp_k\right)t_P+
                   \left(\int_Yz(p_j-p_k)\right)t_\Sigma.
 \end{aligned}
\end{equation}
The relative middle matrix is
$B'=\bigl(\begin{smallmatrix}3&-2\\-2&4\end{smallmatrix}\bigr)$
from \eqref{eq:ruled-matrix}. Its cokernel map is
$(n_P,n_\Sigma)\mapsto6n_P+n_\Sigma\pmod8$.
The boundary has the groups \eqref{eq:L-groups}.
Write $\tau_k$ for the image of $t_\Sigma$ and retain the
symbols $\xi,\upsilon$ for their boundary images. Its nonzero
positive products are
\begin{equation}\label{eq:ruled-L}
 \xi^2=6\tau_k,\quad \xi\upsilon=\upsilon^2=\tau_k,
 \quad \xi \eta_\xi=\upsilon \eta_\upsilon=\omega,\quad8\tau_k=0,
\end{equation}
where $\eta_\xi,\eta_\upsilon$ are the dual degree-five basis.
Their cross pairings are zero. Every other positive product
vanishes.

\begin{proposition}\label{prop:all-attachments}
In every ruled row of \eqref{eq:ruled-template}, restriction
from the complement is
\begin{equation}\label{eq:ruled-rho}
 \rho^2(a)=s\xi+t\upsilon,\qquad
 \rho^4([z])=\left(\int_Yz(p_j+5p_k)\right)\tau_k\pmod8.
\end{equation}
Its relative maps in degrees $4,6,8$ are
\begin{equation}\label{eq:ruled-relative}
 (p_k,p_j-p_k),\qquad
 (\epsilon\ell_k,-\epsilon(\ell_0+\ell_2)),\qquad[\pt].
\end{equation}
For the ordered difference $r_V-r_W$, the connecting maps
on $\eta_\xi,\eta_\upsilon,\omega$ are the negatives of the
corresponding line and point maps. All six complement
rings are $R/(p_0,p_2)$.
\end{proposition}

\begin{proof}
Proposition~\ref{prop:compact-walls} and
\eqref{eq:surface-transform} give these as the actual
adjacent compact configurations. The two surface classes
generate the same ideal $(p_0,p_2)$, and their line pair
is a unimodular change of $\ell_0,\ell_2$.
The relative groups and primitive maps therefore give
the same quotient proof as in
Proposition~\ref{prop:complement}. Their regular
neighborhoods have the computed component ring and
intersection matrix, so Poincar\'e--Lefschetz duality
gives \eqref{eq:ruled-L} and \eqref{eq:ruled-rho}.
The actual compact maps identify the exteriors, and
regular levels outside the exceptional core identify
their outward boundaries. Thus these are the overlap
maps of the given geometric gluing.

To check integrality and multiplicativity explicitly,
evaluation of $p_j+5p_k$ against $p_j,p_k$ gives $8,16$,
so the degree-four map kills the ideal. For $a,b\in R^2$,
\[
 \int_Yabp_k=d_k(a)d_k(b),\qquad
 \int_Yab(p_j-p_k)=d_j(a)d_j(b)-d_k(a)d_k(b).
\]
Substitution of $s,t$ in \eqref{eq:ruled-L} gives
$6s_as_b+s_at_b+t_as_b+t_at_b$, which is the same
value modulo eight. All other positive products have
zero target. The relative maps are the geometric
component, line and point inclusions. Their connecting
signs follow from the same outward Stokes convention as
\eqref{eq:MV}.
\end{proof}

For $k=2$ the functional in \eqref{eq:ruled-rho} is
$\kappa$, while for $k=0$ it is $5\kappa$. This change
records the torsion-generator convention, not a new
uncomputed gluing choice. For example across $A\to B$,
the degree-two boundary classes transform as
\[
 x_0\longmapsto\upsilon,\qquad
 x_2\longmapsto\upsilon-\xi.
\]
The degree-five classes transform by the inverse transpose:
$\eta_0\mapsto \eta_\xi+\eta_\upsilon$, $\eta_2\mapsto-\eta_\xi$,
and $\omega\mapsto\omega$. These formulas preserve
\eqref{eq:L-products} and \eqref{eq:ruled-L}, with
$\tau\mapsto\tau_2$.

\subsection{The moving divisor remains in the complement}
The open set $V_\sigma$ contains the divisor $D^\circ$
over $S^\circ$, so it is larger than $\pi_\sigma^{-1}(U)$.
We record its established ring to make this distinction
explicit. In chamber $A$ use $x$ from \eqref{eq:witness},
set $z_D=a_1+2x$, and put
\[
 d=\theta_w(z_D)|_{S^\circ},\qquad
 h_D=\theta_v(x)|_{D^\circ}.
\]
The open-divisor calculation of \cite[\S7]{wuebben17} is
\begin{equation}\label{eq:moving-ring}
 \begin{gathered}
 H^2(S^\circ,\Z)=H^2(S,\Z)/\Z[C]\simeq\Z^{21},\qquad
 H^k(S^\circ,\Z)=0\quad(k>0,\ k\ne2),\\
 H^*(D^\circ,\Z)=H^*(S^\circ,\Z)[h_D]/(h_D^2-dh_D),\qquad
 D|_{D^\circ}=-2h_D+\pi^*d.
 \end{gathered}
\end{equation}
Here $[C]$ is primitive, so the Gysin sequence gives the
base cohomology. The $\mathrm{PGL}_2$ structure group acts
trivially on fibre cohomology, and $h_D$ has fibre degree
one, giving an integral Leray--Hirsch basis. The topological
rank-two lift exists because $H^3(S^\circ,\Z)=0$;
normalizing its hyperplane class to $h_D$ gives first
Chern class $d$ and zero second Chern class. The universal
extension gives the divisor identity; the orthogonality
of $z_D$ to $a_1,w,v$ removes base-line ambiguities.

The algebraic Brauer obstruction remains nonzero of order
two on $S^\circ$: restriction to the function field is
injective for a regular integral surface. Its analytic
image on this open surface vanishes, using
$H^2(S^\circ,\OO)=H^3(S^\circ,\Z)=0$. These lifting
statements and their primary sources are proved in
\cite[\S7.1]{wuebben17}. They do not imply an algebraic
universal family over the full compact $S$. The divisorial
wall proof used coarse blow-ups precisely so that such a
lift was unnecessary.
